% tex/frontmatter/abstract.tex
\chapter*{Abstract}
\addcontentsline{toc}{chapter}{Abstract}
\markboth{Abstract}{Abstract}

% TODO: Write abstract (max ~300 words, no figures/tables/equations/references).
%
% The abstract should:
%   (1) state the scope and principal objectives of the research,
%   (2) describe the methods used,
%   (3) summarise the results, and
%   (4) state the principal conclusions.

Avalanches are the most fatal natural hazard in Switzerland \parencite{LongtermStatistics}, and while processes like crack propagation, the transition between crack propagation regimes and the influence of snow and terrain parameters have been studied, there is little link between crack propagation and its seismic signature. In this master thesis, spectral element modelling in SALVUS will be used to simulate crack propagation in snow, as measured by distributed acoustic sensing (DAS). The spectral element modelling results will encompass simple simulations of sub-rayleigh and supershear propagation of uniform source wavelets in snow to determine the different types of waves in snow. Furthermore, the results from material point method simulations will be used to obtain seismic moment tensors and source wavelet shapes in an attempt to make the crack propagation model more physical. \par

The simple model results show that the main energy in the results is focused at the propagation speed of the crack. In the simple sub-rayleigh case, there is a precursor to the main energy front, which has been identified to be the S0 wave because of its speed and the plate wave speed, which is consistent with \cite{herwijnenSeismicSensorArray2011}. This precursor is caused by lack of interference of the first source at dense spacing. Overall, sub-rayleigh regimes have lower energy than supershear propagation regimes, which is important because they might not be detected in DAS, especially if there is a transition in propagation speed, as observed by \cite{trottetTransitionSubRayleighAnticrack2022}. Another difference between supershear and sub-rayleigh crack propagation is that supershear wavefronts outside the main rupture path decay in amplitude with increasing time in the space-time plots, whereas sub-rayleigh have a constant amplitude. Both regimes show a rupture arrest front, which occurs in different modes. In space-time, supershear cracks have a fan-like front which only extends in the positive direction, whereas sub-rayleigh has a v-shaped arrest front which propagates in both the positive and negative direction. This is important for identifying the propagation regime in DAS field data. \par 

The material-point method guided results further support the effects observed in the simple model. The main wavefront, which travels at sub-rayleigh speed, shows similar signatures. The S0-precursor in the space-time plots can also be observed. The same is true for the v-shaped rupture arrest signatures in space-time, which were observed in the simple sub-rayleigh model. Furthermore, the MPM-guided results show a quasi-static strain field, which is absent from the simple model. This quasi-static strain field is a precursor to the main wavefront and could be used for early warning of avalanche nucleation. The amplitude of this quasi-static strain field is higher than the S0 precursor, which means that it might be more easily detected in DAS field data. \par

\cleardoublepage
